\documentclass[../../main.tex]{subfiles}
\graphicspath{{../../image/}{../../figure/}}
\begin{document}

\section{真正非线性标量平衡律}\label{section:ch11-真正非线性标量平衡律}

现转向守恒律\eqref{align:ch11-11-1-1}在带\textbf{源项}情形的大时间行为, 即
\begin{align}
\partial_t u(x,t)+\partial_x f(u(x,t))=g(u(x,t),x,t).\label{align:ch11-11-9-1}
\end{align}
线性源项 $g=au$ 的情形将在下节单独处理。此处假设 $f''(u)>0$, $-\infty<u<\infty$, 即与本章前八节相同的真正非线性条件。

为阐明源项的影响, 先回忆在右端为零的情形, 所有激波都以通过间断两端的特征线加权平均的速度传播, 并且任意凸熵的熵产生在激波集上消失。特别地, 从 $(x,t)$ 出发的极小与极大后向特征线, 斜率为 $f'(u(x-,t))$ 与 $f'(u(x+,t))$, 夹住激波轨迹, 速度为中间值 $[f(u(x+,t))-f(u(x-,t))]/[u(x+,t)-u(x-,t)]$。

在右端非零情形, 我们仍可定义一个 Lipschitz 通量
\begin{align}
j(u)=f(u)-f(0)-f'(0)u,\label{align:ch11-11-9-2}
\end{align}
它生成一族凸熵 $\eta_v(u)=(u-v)^+$, 熵流
\begin{align}
q_v(u)=
\begin{cases}
f(u)-f(v), & u>v,\\
0, & u<v.
\end{cases}\label{align:ch11-11-9-3}
\end{align}
对每个固定 $v$, 可与解 $u$ 关联熵产生
\begin{align}
E_v=\partial_t(u-v)^++\partial_x q_v(u)-g(u,x,t)\chi_{\{u>v\}}.\label{align:ch11-11-9-4}
\end{align}

\begin{theorem}\label{theorem:ch11-11-9-1}
存在与可容许解 $u$ 关联的广义特征线 $\xi(\cdot)$, 在时间区间 $[\sigma,\tau]\subset(0,\infty)$ 上满足
\begin{align}
\dot{\xi}(t)\in
\begin{cases}
\{f'(u(\xi(t)\pm,t))\}, & \text{当 }u(\xi(t)+,t)=u(\xi(t)-,t),\\[1mm]
\left[\dfrac{f(u(\xi(t)+,t))-f(u(\xi(t)-,t))}{u(\xi(t)+,t)-u(\xi(t)-,t)},f'(u(\xi(t)-,t))\right], & \text{当 }u(\xi(t)+,t)<u(\xi(t)-,t),
\end{cases}\label{align:ch11-11-9-5}
\end{align}
对几乎所有 $t\in[\sigma,\tau]$。从任意点 $(\bar x,\bar t)\in(-\infty,\infty)\times(0,\infty)$ 出发, 存在唯一极小与唯一极大后向广义特征线, 进入区域 $\{t<\bar t\}$。
\end{theorem}

\begin{proof}
定义广义特征线, 例如极小者, 为满足\eqref{align:ch11-11-9-5}中下半闭特征锥的绝对连续曲线 $\xi(\cdot)$。存在性可由标准逼近论证, 用离散网格逼近上半平面而得; 极小性和唯一性由闭锥性质推出。激波两侧特征线之间的楔形区域由熵条件 $u_-\geqslant u_+$ 推出。
\end{proof}

应当强调, 与源项 $g$ 的大小无关, 特征锥由通量 $f$ 完全确定。因此, 上一节发展的分界线理论仍适用: 分界线, 即从 $(0,0)$ 发出、无激波的广义特征线, 是 $u$ 为常数的 Lipschitz 曲线。从 $x$ 轴上的任意点出发, 总可构造至少一条前向特征线, 它是初始点与其特征锥之间广义特征线的极大者。这一理论特别推出:

\begin{theorem}\label{theorem:ch11-11-9-2}
假设初值 $u_0$ 在 $(-\infty,\infty)$ 上可积, 并设
\begin{align}
q_-=\max_x\int_x^{-\infty}u_0(y)\,\mathrm{d}y,\qquad q_+=\max_x\int_x^\infty u_0(y)\,\mathrm{d}y.\label{align:ch11-11-9-6}
\end{align}
则存在 Lipschitz 曲线 $\chi_-(\cdot)$ 与 $\chi_+(\cdot)$, 分界线, 使得
\begin{align}
\max_x\int_x^{-\infty}u(y,t)\,\mathrm{d}y=q_-,\qquad
\max_x\int_x^\infty u(y,t)\,\mathrm{d}y=q_+,\label{align:ch11-11-9-7}
\end{align}
并在 $\chi_-(t)$ 与 $\chi_+(t)$ 处取到。
\end{theorem}

\begin{proof}
分界线, 即从使\eqref{align:ch11-11-9-6}中极大取到的点 $(x_\pm,0)$ 出发的广义特征线, 是 Lipschitz 曲线, 与零源情形一样携带质量 $q_\pm$, 因为它们夹住在其上 $u=0$ 的区域。
\end{proof}

注意\eqref{align:ch11-11-9-7}是零源情形\cref{theorem:ch11-11-4-2}的精确类似, 但我们现在利用了 $j(0)=0=j'(0)$。由熵条件 $u_-\geqslant u_+$, 激波速度落在斜率 $f'(u_-)$ 与 $[f(u_+)-f(u_-)]/[u_+-u_-]$ 之间, 而非后者与 $f'(u_+)$ 之间。因此, 楔形不对称地位于激波两侧, 反映源项 $g$ 引起的激波运动修正。

此框架下还有一条性质保持: 激波消失的广义特征线是经典特征线。

\begin{theorem}\label{theorem:ch11-11-9-3}
设 $\xi(\cdot)$ 是与可容许解 $u$ 关联、在时间区间 $[\sigma,\tau]$ 上无激波的广义特征线。则 $u$ 在 $\xi(\cdot)$ 上满足
\begin{align}
\frac{d}{dt}u(\xi(t),t)=g(u(\xi(t),t),\xi(t),t).\label{align:ch11-11-9-8}
\end{align}
特别地, $u(\xi(t)\pm,t)$ 是 $t$ 的 Lipschitz 函数。
\end{theorem}

\begin{proof}
在区域 $\{(x,s):\sigma<s<t,\ \xi(s)-\varepsilon<x<\xi(s)\}$ 与 $\{(x,s):\sigma<s<t,\ \xi(s)<x<\xi(s)+\varepsilon\}$ 上分别积分守恒律, 应用 Green 公式, 并令 $\varepsilon\downarrow0$。由于 $u$ 在 $\xi(\cdot)$ 上连续, \eqref{align:ch11-11-9-5}给出 $\dot{\xi}=f'(u)$, 沿该曲线得到\eqref{align:ch11-11-9-8}。
\end{proof}

由此, 与 $u$ 关联、在 $[\sigma,\tau]$ 上无激波的两条不同广义特征线, 在 $(\sigma,\tau)$ 内不相交。特别地, 极值后向特征线可分开逼近任意间断点。

现转至熵产生, 它不再只集中在激波集上。

\begin{theorem}\label{theorem:ch11-11-9-4}
对任意 $v\in\mathbb{R}$, \eqref{align:ch11-11-9-4}定义的熵产生 $E_v$ 是上半平面上的非负 Radon 测度。此外, $E_v$ 在 $u(\cdot,t)$ 的连续点集 $\mathscr{C}$ 的内部消失:
\begin{align}
E_v\llcorner\mathscr{C}^0=0.\label{align:ch11-11-9-9}
\end{align}
\end{theorem}

\begin{proof}
非负性是可容许性的推论, 即所有熵不等式 $\partial_t\eta(u)+\partial_x q(u)\leqslant0$(在零源情形)以及它们在源项存在时的修正版本。为在区域 $\mathscr{C}^0$ 上建立消失, 固定 $(\bar x,\bar t)\in\mathscr{C}^0$。设 $\xi(\cdot)$ 与 $\zeta(\cdot)$ 是从 $(\bar x,\bar t)$ 出发的极小与极大后向特征线。由\eqref{align:ch11-11-9-8}, $u$ 沿二者均满足常微分方程, 且因 $(\bar x,\bar t)$ 是连续点, 故 $\xi=\zeta$。因此, 三角形区域 $S=\{(x,t):0<t<\bar t,\ \xi(t)<x<\zeta(t)\}$ 收缩为曲线, 在其上积分\eqref{align:ch11-11-9-4}给出
\begin{align}
\int_S E_v=\lim_{\varepsilon\to0}\int_0^{\bar t}\int_{\xi(t)-\varepsilon}^{\zeta(t)+\varepsilon}\partial_t\eta_v+\partial_x q_v-g\chi_{\{u>v\}}\,\mathrm{d}x\mathrm{d}t=0,\label{align:ch11-11-9-10}
\end{align}
其中 $\eta_v(u)=(u-v)^+$。

然而, 存在另一可能来源: 相邻激波的碰撞, 这会在不规则点产生熵产生。但由\cref{theorem:ch11-11-3-4}, 不规则点集可数, 故不会贡献绝对连续测度部分。综上,
\begin{align}
E_v=e_v(t)\,\mathrm{d}\mu_{\mathscr{J}},\label{align:ch11-11-9-11}
\end{align}
即 $E_v$ 被 $u$ 的跳跃集 $\mathscr{J}$ 携带。因此在连续点内部 $\mathscr{C}^0$ 上消失。
\end{proof}

\end{document}
